% FULL SAMPLE PAPER — complete text from abstract, introduction, related work,
% methodology, experiments (template.tex), and discussion source files.
% Compile: pdflatex sample_template && bibtex sample_template && pdflatex x2

\documentclass{article}
\usepackage{arxiv}
\usepackage[utf8]{inputenc}
\usepackage[T1]{fontenc}
\usepackage{hyperref}
\usepackage{url}
\usepackage{booktabs}
\usepackage{amsfonts}
\usepackage{microtype}
\usepackage{graphicx}
\usepackage{natbib}
\usepackage{tabularx}
\usepackage{multirow}
\usepackage{diagbox}
\usepackage{array}
\usepackage{amsmath}
\usepackage{algorithm}
\usepackage{algorithmic}
\usepackage{float}
\usepackage{subcaption}

\newcommand{\maybeincludegraphics}[2][]{%
  \IfFileExists{#2}{%
    \includegraphics[#1]{#2}%
  }{%
    \fbox{\parbox[c][2.2cm][c]{0.88\linewidth}{\centering\small\textit{Figure file not found:}\ \texttt{#2}}}%
  }%
}

\title{DIVERGE with Semantic Structure Enhancement: Pruning Noisy Edges for Text-Attributed Graph Learning}

\author{
  Nima Parifard \\
  Department of Computer Science \\
  AmirKabir University \\
  \texttt{nima.parifard@aut.ac.ir}
  \And
  Saeedeh Momtazi \\
  Department of Computer Science \\
  AmirKabir University \\
  \texttt{momtazi@aut.ac.ir}
}

\renewcommand{\shorttitle}{DIVERGE + Semantic Structure Enhancement}

\begin{document}
\maketitle

\begin{abstract}
Node classification on text-attributed graphs (TAGs) requires both rich textual
semantics and a graph topology aligned with the task. Large language model (LLM)
pipelines often improve performance by \emph{rewriting} node text or structure, but
such enhancers can be costly, brittle, and dependent on external API access. The
\textbf{DIVERGE} framework takes a different path: it fine-tunes a pre-trained LLM with
LoRA under multiple random initializations to obtain diverse encodings of \emph{native}
node text, trains one graph neural network (GNN) per view, and combines their outputs
with adaptive ensemble learning rather than uniform logit averaging. We show that even
strong LLM representations are limited when observed edges connect semantically
dissimilar nodes, violating homophily and degrading message passing.

We extend DIVERGE with a \textbf{semantic structure enhancer}---a lightweight,
post-hoc module that prunes noisy edges before per-view GNN training, without modifying
node text or calling LLMs for topology editing. Two criteria are proposed:
(1)~\emph{cosine-similarity pruning}, which removes edges whose LLM endpoint embeddings
fall below a validation-tuned threshold; and
(2)~\emph{uncertainty-based pruning}, which removes edges with highest predictive
uncertainty from a cross-validated GNN edge predictor (standard deviation or entropy).
Each LLM view yields its own refined graph, so the ensemble jointly exploits semantic,
structural, and predictive diversity.

Experiments on Cora, CiteSeer, PubMed, and WikiCS show that targeted edge removal
consistently improves over DIVERGE alone and outperforms strong baselines including
TAPE, while random edge deletion severely hurts accuracy---confirming that gains arise
from informed denoising rather than sparsification alone. Cosine and uncertainty
enhancers are complementary across datasets, offering an efficient alternative to
full LLM enhancer pipelines for TAG learning.
\end{abstract}

\keywords{Text-Attributed Graphs \and Large Language Models \and Graph Neural Networks \and Graph Structure Learning \and Ensemble Learning}

%=============================================================================
\section{Introduction}
\label{sec:introduction}
%=============================================================================

Text-attributed graphs (TAGs) are ubiquitous in real-world systems: citation networks
link scholarly documents, e-commerce graphs connect products through co-purchase or
review relations, and social platforms relate user-generated content through follow or
interaction edges \citep{Zhang2024TAGSurvey}. In a TAG, each node is associated with a
natural-language description and each edge encodes a relational tie. Node classification
---predicting the category of each node from its text and neighborhood---is a fundamental
task that requires jointly exploiting \emph{semantic content} and \emph{graph structure}.
Graph neural networks (GNNs) excel at relational reasoning through message passing, but
their effectiveness depends critically on the quality of node features and on whether
the observed topology aligns with the homophily assumption that connected nodes tend to
be similar \citep{Kipf2016GCN,Velickovic2018GAT,jin2020prognn}.

Recent progress has shifted TAG pipelines from shallow text encoders toward large
language models (LLMs), which capture richer semantics from node documents
\citep{Chen2024LLaGA,Huang2024GNNAdapters,he2024harnessing}. A parallel research line
treats LLMs as \emph{enhancers}: they rewrite node text, relabel nodes, synthesize new
vertices, or modify graph structure before GNN training \citep{he2024harnessing,
Sun2026TopologicalEnhancerTAG,zhang2025leveraging,Yu2025NodeGenerationTAG}. Although
powerful, enhancer methods often incur API cost, hallucination risk, or multi-stage
complexity, and they alter the native textual content of the graph. An alternative
strategy asks whether strong performance can arise from \emph{diverse encodings of
unchanged node text} rather than from generative augmentation.

\textbf{DIVERGE} addresses this question by fine-tuning a pre-trained LLM with
parameter-efficient adaptation (LoRA) under multiple random initializations. Each run
produces a distinct semantic view of the same native documents; a separate GNN is
trained on each view, and an adaptive ensemble combines their logits---assigning
learnable weights at the model or class level rather than uniformly averaging predictions
as in TAPE \citep{he2024harnessing}. DIVERGE thus pursues diversity in \emph{how} text
is encoded, not in \emph{what} text is provided to the model.

Despite strong semantic representations, a second challenge remains largely unaddressed
in this setting: \textbf{the observed graph structure may be noisy or misaligned with
semantics}. Citation links, hyperlinks, and co-occurrence edges are constructed by
rules that do not guarantee semantic similarity between endpoints. When dissimilar nodes
are connected, message passing propagates irrelevant information, lowers effective
homophily, and limits GNN performance even if node embeddings are high quality
\citep{zhu2021survey,chen2020idgl,dong2023gps}. Classical graph structure learning (GSL)
methods such as ProGNN, IDGL, and PTDNet refine topology using node features
\citep{jin2020prognn,chen2020idgl,luo2021ptdnet}, but they typically operate on a single
shallow feature space in a joint end-to-end pipeline. LLM-based structural enhancers,
meanwhile, rewrite edges or relation types through external model calls
\citep{Sun2026TopologicalEnhancerTAG,guo2024graphedit,zhang2025leveraging}. We argue
that DIVERGE's multi-view LLM embeddings offer a natural signal for \emph{lightweight,
post-hoc} structure refinement---without generating new text or invoking additional LLM
queries at enhancement time.

\paragraph{Our approach.}
We extend DIVERGE with a \textbf{semantic structure enhancer} that prunes noisy edges
before per-view GNN training and ensemble aggregation. The key insight is that nodes
with similar semantics cluster in LLM embedding space; edges whose endpoints have low
cosine similarity are likely to harm homophily-driven message passing. We propose two
complementary mechanisms:

\begin{itemize}
\item \textbf{Cosine-similarity enhancement (Method~II).} For each LLM view
$\bar{f}_i$, we compute node embeddings $\mathbf{Z}^{(i)}$ and remove edges
$(v_i,v_j)$ whose endpoint representations satisfy
$\cos(\mathbf{z}_i,\mathbf{z}_j) < T$ for a threshold~$T$. This yields a collection of
view-specific refined graphs $\{\bar{G}_i\}_{i=1}^{M}$, each with higher semantic
homophily while preserving native text and node features.

\item \textbf{Uncertainty-based enhancement (Method~III).} We train a GNN-based edge
predictor with cross-validation and treat edges with high predictive uncertainty---measured
by entropy or standard deviation across folds---as unreliable. Removing the most
uncertain edges again increases homophily and complements similarity-based pruning with
a model-driven criterion related to edge-prediction GSL
\citep{yu2020grcn,zhao2021gaug}.
\end{itemize}

Both enhancers plug into the existing DIVERGE pipeline: multiple LoRA fine-tunes,
per-view GNN training on refined graphs, and adaptive ensemble learning over logits.
The final model therefore exploits \emph{three} sources of diversity---semantic
(repeated LLM encodings), structural (view-specific refined topologies), and predictive
(learned ensemble weights).

\paragraph{Contributions.}
Our main contributions are summarized as follows:
\begin{itemize}
\item We identify a structural bottleneck in LLM-on-TAG pipelines: even diverse,
  high-quality semantic embeddings cannot fully compensate for edges that connect
  semantically dissimilar nodes, and we motivate homophily-oriented edge pruning in
  this setting.

\item We propose two semantic structure enhancers for DIVERGE---cosine-similarity
  thresholding and uncertainty-based edge removal---that refine topology using LLM
  embeddings or GNN edge predictions \emph{without} rewriting node text or calling
  external LLM APIs for structure modification.

\item We integrate structural refinement with DIVERGE's multi-view encoding and
  adaptive ensemble framework, producing view-specific denoised graphs rather than a
  single jointly learned adjacency as in classical GSL
  \citep{chen2020idgl,jin2020prognn}.

\item We conduct extensive experiments on benchmark TAG datasets, showing that
  removing noisy edges consistently improves both individual GNN accuracy and ensemble
  performance compared to DIVERGE alone and to representative LLM-enhancer baselines
  such as TAPE \citep{he2024harnessing}.
\end{itemize}

\paragraph{Paper organization.}
The remainder of this paper is organized as follows.
Section~\ref{sec:related-works} reviews graph structure learning, TAG-specific
denoising, and LLM enhancer methods.
Section~\ref{sec:methodology} presents notation, a brief DIVERGE recap, and the
cosine-similarity and uncertainty-based structural enhancers.
Section~\ref{sec:experiments} reports experimental results and ablations.
Sections~\ref{sec:discussion} and~\ref{sec:conclusion} discuss findings, limitations,
and future directions.


%=============================================================================
\section{Related Works}
\label{sec:related-works}
%=============================================================================

Text-attributed graphs (TAGs) combine relational structure with natural-language
node descriptions, making them well suited to citation networks, product graphs,
and social platforms where both topology and text carry predictive signal
\citep{Zhang2024TAGSurvey}. A central design question is how to fuse textual
semantics with graph propagation: early pipelines encoded documents with
shallow features (e.g., bag-of-words) and fed the resulting vectors into graph
neural networks (GNNs), whereas recent work leverages pre-trained language models
(PLMs/LLMs) for richer node representations \citep{Chen2024LLaGA,Huang2024GNNAdapters}.
Our work builds on \textbf{DIVERGE}, which fine-tunes an LLM with parameter-efficient
adaptation (LoRA) under \emph{multiple random initializations} to obtain diverse
semantic encodings of \emph{native} node text---without rewriting, paraphrasing,
or augmenting the original documents---and then trains one GNN per encoding before
combining their outputs with adaptive ensemble learning rather than uniform logit
averaging \citep{he2024harnessing}. In this paper we extend DIVERGE with
\textbf{semantic structure enhancement}: we refine the observed topology when
edges connect semantically dissimilar nodes, thereby increasing homophily before
message passing. We organize prior work into three lines---general graph structure
learning, structure-aware TAG learning without LLM enhancers, and LLM-based
enhancement---and then state the research gap our approach addresses.


% ---------------------------------------------------------------------------
\subsection{Edge Rewiring and Denoising in Graph Learning}
\label{sec:gsl-general}
% ---------------------------------------------------------------------------

Real-world graphs are often noisy, incomplete, or misaligned with node features.
Because GNNs propagate information along observed edges, spurious connections can
inject irrelevant signals and degrade node classification
\citep{zhu2021survey,jin2020prognn}. \emph{Graph structure learning} (GSL) addresses
this by jointly or sequentially refining topology and GNN parameters.

\paragraph{Joint structure--parameter optimization.}
\textbf{LDS} \citep{franceschi2019lds} formulates structure learning as bilevel
optimization over a discrete edge distribution and GCN weights, enabling learning
when graphs are missing or corrupted. \textbf{ProGNN} \citep{jin2020prognn} treats
the adjacency matrix as a learnable variable regularized toward sparsity, low rank,
and feature smoothness, and has become a standard reference for graph denoising under
adversarial or noisy edges. \textbf{IDGL} \citep{chen2020idgl} iteratively alternates
between (i) constructing a new graph from \emph{cosine similarity} of node embeddings
and (ii) retraining the GNN on the updated graph until convergence. This iterative,
similarity-driven rewiring is particularly relevant to our cosine-based structural
enhancer (Method~II), although IDGL operates end-to-end on shallow features rather
than on post-hoc LLM embeddings from multiple semantic views.

\paragraph{Explicit edge removal and feature-aware rewiring.}
Several methods target \emph{denoising} rather than full adjacency re-estimation.
\textbf{PTDNet} \citep{luo2021ptdnet} learns a parameterized edge-drop module trained
jointly with the downstream GNN to remove task-irrelevant connections.
\textbf{GPS} \citep{dong2023gps} introduces an edge signal-to-noise ratio (ESNR) to
quantify structural noise and rewires graphs using self-supervised link propensity
scores derived from node features. \textbf{DIGL} \citep{klicpera2019digl} smooths
the adjacency via personalized PageRank diffusion before propagation, offering a
classic pre-processing denoising baseline. Together, these works establish that
\emph{pruning or rewiring unreliable edges} can substantially improve GNN performance
when node features encode meaningful similarity---a motivation we adopt in the
TAG setting using LLM-derived semantics.

\paragraph{Edge prediction and graph revision.}
\textbf{GRCN} \citep{yu2020grcn} couples a graph-revision GNN that adds or reweights
edges with a node-classification GNN optimized jointly on the revised structure.
\textbf{GAug} \citep{zhao2021gaug} uses a graph autoencoder to predict edge existence,
adding likely links and dropping unlikely ones before downstream training. These
edge-predictor designs parallel the spirit of our uncertainty-based enhancer
(Method~III), where an encoder--decoder estimates edge reliability; however, GRCN
and GAug revise structure inside a single end-to-end pipeline on fixed input features,
whereas we use cross-validated prediction uncertainty as a separate, interpretable
criterion for noisy-edge removal on graphs already enriched by diverse LLM encodings.

\paragraph{Self-supervised and compact structure learning.}
\textbf{STABLE} \citep{li2022stable} first learns robust node embeddings via
contrastive self-supervision and then rewires the graph by embedding similarity.
\textbf{CoGSL} \citep{fatemi2022cogsl} compresses mutual information across multiple
graph views to obtain a minimal sufficient structure for the task.
\textbf{SUBLIME} \citep{liu2022sublime} performs unsupervised GSL through structure
bootstrapping and contrastive agreement between learned and anchor graphs. These
methods demonstrate that structure can be improved even with limited or no direct
supervision on edges, complementing our homophily-driven pruning strategies.


% ---------------------------------------------------------------------------
\subsection{Edge Denoising for Text-Attributed Graphs}
\label{sec:tag-structure-no-llm}
% ---------------------------------------------------------------------------

On TAGs, structure quality is equally critical: citation, co-purchase, and hyperlink
graphs frequently connect documents that are topologically related yet semantically
distant, violating the homophily assumptions underlying standard message passing
\citep{yang2019vhe,Zhang2024TAGSurvey}. Early textual-network models such as
\textbf{VHE} \citep{yang2019vhe} encode homophily directly in a variational
framework, encouraging correlated latent representations for linked nodes while
jointly modeling text and topology; however, VHE does not explicitly remove noisy
edges from the observed graph. Joint graph--GNN learners including \textbf{GLCN}
\citep{jiang2019glcn} and \textbf{GEN} \citep{wang2021gen} estimate or refine
adjacency from node features and have been widely evaluated on citation-network
benchmarks (Cora, Citeseer, PubMed) with bag-of-words attributes---historically the
standard TAG setting before LLM encoders became dominant.

Importantly, most GSL methods above were \emph{not} designed for modern TAG pipelines
where node features come from fine-tuned LLMs and where multiple semantic views of
the same native text are available. They typically (i) use a single feature source,
(ii) entangle structure learning with one GNN training run, and (iii) do not exploit
representation diversity from repeated parameter-efficient fine-tuning. Our structural
enhancers differ precisely along these axes: we keep native text unchanged, derive
\emph{multiple} LLM embedding spaces from DIVERGE, and apply \emph{post-hoc}
topology refinement---cosine-threshold pruning (Method~II) or uncertainty-based edge
removal (Method~III)---before training the per-view GNNs and ensemble. This positions
our work between classical GSL on citation graphs and recent LLM enhancers that modify
text or structure through external model calls (next subsection).


% ---------------------------------------------------------------------------
\subsection{LLMs as Enhancers for Text-Attributed Graph Learning}
\label{sec:llm-enhancers}
% ---------------------------------------------------------------------------

A growing line of work assigns LLMs an \emph{enhancer} role: they alter node text,
graph structure, or labels \emph{before or during} GNN training to improve downstream
performance \citep{Jin2023LLMGraphs,Chen2024LLMGraphs,Li2024GraphLLMSurvey,Wu2025LLMNode}.
We group enhancer methods by the modality they modify.

\paragraph{Text and feature enhancement.}
\textbf{TAPE} \citep{he2024harnessing} queries an LLM for zero-shot predictions and
natural-language explanations, translates explanations into dense features through an
LLM-to-LM interpreter, and trains multiple GNNs whose logits are uniformly averaged.
Although highly effective, TAPE depends on external LLM API access and may inject
semantic noise when native texts are already informative. \textbf{NRUP}
\citep{kuang2023unleashing} constructs a hierarchical text--word graph and pre-trains
with self-supervised tasks that co-update text and structure. \textbf{E-LLaGNN}
\citep{Jaiswal2024AllAgainstSome} limits cost by enhancing only a budgeted subset of
critical nodes via on-demand LLM prompts, enabling LLM-free inference at deployment.
These methods enrich \emph{inputs} to the GNN; DIVERGE instead seeks gains from
\emph{diverse encodings of unchanged native text} combined with adaptive ensembling.

\paragraph{Structure and topology enhancement.}
Structural LLM enhancers modify edges or relation types rather than rewriting node
documents. \textbf{RoSE} decomposes single-relation edges into multiple semantic
relation types using LLM-based generation and classification of edge semantics.
\textbf{LLM4GraphTopology} \citep{Sun2026TopologicalEnhancerTAG} uses LLM-estimated
semantic similarity to add or delete edges and applies pseudo-label propagation as
topology regularization during GNN training. \textbf{GraphEdit} \citep{guo2024graphedit}
instruction-tunes an LLM to denoise and complete graph structure from a global
reasoning perspective. \textbf{Locle} \citep{zhang2025leveraging} performs active
label-free learning with LLM pseudo-labels and includes GNN-guided topology rewiring
under a fixed query budget. Like our work, these methods acknowledge that observed TAG
topology may be suboptimal; unlike our work, they rely on LLM rewriting of structure
or labels rather than scoring edge reliability from DIVERGE's multi-view semantic
embeddings and lightweight GNN-based uncertainty.

\paragraph{Data generation, label efficiency, and unified pipelines.}
\textbf{LLM4NG} \citep{Yu2025NodeGenerationTAG} synthesizes labeled nodes and predicted
edges for few-shot TAG learning. \textbf{GAGA} \citep{hu2025gaga} annotates a small
representative subset of nodes/edges with an LLM and aligns an annotation graph to the
full TAG. \textbf{UltraTAG} \citep{yu2025ultratag} modularizes augmentation, encoding,
and GNN training into a unified robust pipeline. \textbf{SKETCH} \citep{zhou-etal-2025sketch}
decouples semantic and structural aggregation via retrieval over node texts and distant
structural contexts. \textbf{TAGA} \citep{zhang2025taga} builds dual Text-of-Graph and
Graph-of-Text views and aligns them through self-supervised learning. Collectively,
enhancer methods show that LLMs can improve TAG learning by enriching inputs---often
at the cost of API fees, hallucination risk, or complex multi-stage pipelines.


% ---------------------------------------------------------------------------
\subsection{Research Gap}
\label{sec:research-gap}
% ---------------------------------------------------------------------------

Existing work leaves three gaps that motivate our DIVERGE extensions with semantic
structure enhancement.

\paragraph{Gap 1: Structure refinement without LLM rewriting.}
GNN-only GSL methods (ProGNN, IDGL, PTDNet, GPS, GRCN, etc.) demonstrate that edge
denoising and homophily enhancement benefit message passing, but they generally operate
on a \emph{single} shallow feature space and jointly optimize one structure--GNN pair.
LLM enhancers improve TAG inputs by \emph{generating} new text, labels, or topology
via costly external models. DIVERGE already shows that competitive performance is
attainable from \emph{native text alone} when multiple LoRA fine-tunes provide diverse
semantic views and adaptive ensembling replaces uniform combination
\citep{he2024harnessing}. We ask whether comparable structural gains can be achieved
by \emph{refining}---not regenerating---the graph using those same diverse LLM
embeddings, without additional LLM API calls at enhancement time.

\paragraph{Gap 2: Multi-view semantics for edge scoring.}
Prior cosine-similarity GSL (notably IDGL) and feature-aware rewiring (GPS, STABLE)
use one embedding source. DIVERGE produces $M$ representation spaces
$\{\mathbf{Z}^{(i)}\}_{i=1}^{M}$ from independently initialized LoRA runs, yielding
$M$ structurally refined graphs $\{\bar{G}_i\}$ rather than a single denoised topology.
Method~II removes edges whose endpoint embeddings fall below a cosine-similarity
threshold in each view; Method~III removes edges with highest predictive uncertainty
from a cross-validated GNN edge predictor. To our knowledge, combining \emph{multi-view
LLM encoding diversity} with \emph{homophily-oriented edge pruning} in a post-hoc
enhancer has not been systematically studied for TAG node classification.

\paragraph{Gap 3: Complementarity of similarity-based and uncertainty-based denoising.}
Cosine pruning (Method~II) is interpretable and directly encodes semantic homophily in
LLM space, but relies on a fixed similarity threshold. Uncertainty pruning (Method~III)
is model-driven and captures edges for which relational evidence is ambiguous across
data splits, aligning with edge-prediction GSL \citep{yu2020grcn,zhao2021gaug} yet
applied here as a modular step atop DIVERGE rather than as a fully joint rewrite of
adjacency during a single GNN training run. Our experiments treat both enhancers as
optional modules that plug into the DIVERGE ensemble, allowing semantic representation
diversity and structural homophily to be exploited jointly.

\paragraph{Summary of contribution relative to prior work.}
In contrast to LLM enhancers that augment text or structure through generation
\citep{he2024harnessing,Sun2026TopologicalEnhancerTAG,guo2024graphedit,zhang2025leveraging},
and in contrast to classical GSL that denoises graphs from shallow features in a
single pipeline \citep{jin2020prognn,chen2020idgl,luo2021ptdnet,dong2023gps}, we
propose a lightweight \textbf{semantic structure enhancer} for DIVERGE: native node
text is encoded by multiple fine-tuned LLMs; noisy edges are pruned per view by
cosine similarity or predictive uncertainty; GNNs trained on refined graphs feed an
adaptive ensemble. This design preserves DIVERGE's no-augmentation philosophy while
importing the homophily motivation of graph rewiring literature into the modern
LLM-on-TAG setting.


%=============================================================================
\section{Methodology}
\label{sec:methodology}
%=============================================================================

\subsection{Problem Setup}
We study node classification on a text-attributed graph
$G=(V,E,D,Y)$, where $V=\{v_1,\ldots,v_N\}$ is the node set,
$E\subseteq V\times V$ the edge set, $D=\{d_v\}$ the native node texts,
and $Y=\{y_v\}$ labels over $C$ classes. The goal is to predict
$\hat{y}_v$ for unlabeled nodes using both $D$ and $E$ without modifying
the original documents.


% ---------------------------------------------------------------------------
\subsection{Background: DIVERGE (Brief Recap)}
\label{sec:diverge-recap}
% ---------------------------------------------------------------------------

\textbf{DIVERGE} \citep{he2024harnessing} is our base framework. It does \emph{not}
augment or rewrite node text; instead it builds diversity from repeated
parameter-efficient fine-tuning of one pre-trained LLM $f(\cdot;\mathbf{W})$ with
LoRA \citep{Hu2021LoRA}:
\begin{equation}
\bar{f}_i(x)=f(x;\mathbf{W}+\mathbf{B}_i\mathbf{A}_i),
\qquad
(\mathbf{A}^{(0)}_i,\mathbf{B}^{(0)}_i)\sim\mathcal{I}_i,
\quad i=1,\ldots,M.
\label{eq:fbar}
\end{equation}
Each run $i$ yields node embeddings $\mathbf{Z}^{(i)}\in\mathbb{R}^{N\times p}$ by
mean-pooling token hidden states of $\bar{f}_i(d_v)$. A GNN is trained on
$(G,\mathbf{Z}^{(i)})$ to produce logits $\mathbf{L}^{(i)}$. Finally, an
\emph{adaptive} ensemble (model-, class-, or per-class MLP weighting) combines
$\{\mathbf{L}^{(i)}\}_{i=1}^{M}$---unlike uniform averaging in TAPE
\citep{he2024harnessing}. Algorithm~\ref{alg:diverge-short} summarizes this pipeline;
our contribution inserts a structure-refinement step \emph{before} each per-view GNN.


\begin{algorithm}[t]
\caption{DIVERGE base pipeline (unchanged text; optional structure refinement)}
\label{alg:diverge-short}
\begin{algorithmic}
\REQUIRE $G=(V,E,D,Y)$, pre-trained LLM $f$, LoRA runs $M$, enhancer $\in\{\emptyset,\ \text{cosine},\ \text{uncertainty}\}$
\ENSURE Predicted labels $\hat{Y}$
\FOR{$i=1$ to $M$}
    \STATE Fine-tune LoRA adapters $\rightarrow$ encoder $\bar{f}_i$ (Eq.~\ref{eq:fbar})
    \STATE Extract embeddings $\mathbf{Z}^{(i)}$ from native texts $D$
    \IF{enhancer $=$ cosine}
        \STATE Build $\bar{G}_i=(V,\bar{E}_i,D,Y)$ by pruning low-similarity edges (Sec.~\ref{sec:cosine-enhancer})
    \ELSIF{enhancer $=$ uncertainty}
        \STATE Build $\bar{G}_i=(V,\bar{E}_i,D,Y)$ by pruning high-uncertainty edges (Sec.~\ref{sec:uncertainty-enhancer})
    \ELSE
        \STATE $\bar{G}_i \leftarrow G$
    \ENDIF
    \STATE Train GNN on $(\bar{G}_i,\mathbf{Z}^{(i)})$; collect logits $\mathbf{L}^{(i)}$
\ENDFOR
\STATE Adaptive ensemble over $\{\mathbf{L}^{(i)}\}$ $\rightarrow$ $\hat{Y}$
\end{algorithmic}
\end{algorithm}


% ---------------------------------------------------------------------------
\subsection{Semantic Structure Enhancement}
\label{sec:structure-enhancer}
% ---------------------------------------------------------------------------

\textbf{Motivation.} LLM embeddings from DIVERGE encode semantic similarity: linked
nodes are often, but not always, close in embedding space. Noisy or rule-induced
edges between dissimilar nodes violate homophily and hurt GNN message passing
\citep{jin2020prognn,chen2020idgl}. We therefore refine $G$ into
$\bar{G}=(V,\bar{E},D,Y)$ with $\bar{E}\subseteq E$ by removing unreliable edges
\emph{before} training each view-specific GNN. Refinement is applied independently
for each embedding $\mathbf{Z}^{(i)}$, yielding
$\mathcal{G}=\{\bar{G}_i\}_{i=1}^{M}$ and models
$\{\mathrm{GNN}_i(\bar{G}_i,\mathbf{Z}^{(i)})\}_{i=1}^{M}$.
Node texts and features are never altered; only topology changes.


% ---------------------------------------------------------------------------
\subsection{Method II: Cosine-Similarity Edge Pruning}
\label{sec:cosine-enhancer}
% ---------------------------------------------------------------------------

For view $i$, let $\mathbf{z}^{(i)}_u$ denote the embedding of node $v_u$ from
$\mathbf{Z}^{(i)}$. For each edge $(v_u,v_v)\in E$, define cosine similarity
\begin{equation}
\mathrm{sim}^{(i)}_{uv}=
\frac{\mathbf{z}^{(i)\top}_u\,\mathbf{z}^{(i)}_v}
{\|\mathbf{z}^{(i)}_u\|_2\,\|\mathbf{z}^{(i)}_v\|_2}.
\label{eq:cosine_similarity}
\end{equation}
Given threshold $T_i\in[-1,1]$, the refined edge set is
\begin{equation}
\bar{E}_i=\{(v_u,v_v)\in E:\ \mathrm{sim}^{(i)}_{uv}\ge T_i\}.
\label{eq:cosine_prune}
\end{equation}
Edges with $\mathrm{sim}^{(i)}_{uv}<T_i$ are treated as semantically inconsistent
(noisy) and removed. Threshold $T_i$ is tuned on validation data (sensitivity analysis
in Section~\ref{sec:experiments}). This is a post-hoc, view-specific analogue of
similarity-based graph structure learning \citep{chen2020idgl} applied to LLM
embeddings rather than jointly learned shallow features.


% ---------------------------------------------------------------------------
\subsection{Method III: Uncertainty-Based Edge Pruning}
\label{sec:uncertainty-enhancer}
% ---------------------------------------------------------------------------

Cosine pruning uses static embedding distance; uncertainty pruning targets edges
whose \emph{existence} is ambiguous under a relational model.

\paragraph{Edge predictor.}
An encoder GNN maps node features (from $\mathbf{Z}^{(i)}$ or initial GNN states
$h_u$) to representations $\mathbf{z}_u=\mathrm{GNN}(h_u)$. A decoder MLP outputs
edge existence probability
\begin{equation}
p_{ij}=\mathrm{MLP}(\mathbf{z}_i,\mathbf{z}_j),
\label{eq:edge_probability}
\end{equation}
trained with binary cross-entropy on observed edges in $E$.

\paragraph{Uncertainty.}
The predictor is retrained on $M_{\mathrm{cv}}$ bootstrap/cross-validation splits,
giving probabilities $\{p_{ij}^{(m)}\}_{m=1}^{M_{\mathrm{cv}}}$ per edge. We use
either standard deviation
\begin{equation}
U_{ij}^{\mathrm{std}}=
\sqrt{\tfrac{1}{M_{\mathrm{cv}}}\textstyle\sum_m\big(p_{ij}^{(m)}-\mu_{ij}\big)^2},
\qquad
\mu_{ij}=\tfrac{1}{M_{\mathrm{cv}}}\textstyle\sum_m p_{ij}^{(m)},
\label{eq:std_uncertainty}
\end{equation}
or entropy $U_{ij}^{\mathrm{ent}}=-\sum_m p_{ij}^{(m)}\log p_{ij}^{(m)}$.
High $U_{ij}$ indicates an unreliable edge.

\paragraph{Pruning strategies.}
Two rules remove noisy edges from $E$:
\begin{itemize}
\item \textbf{Top-$k$\%:} remove the $k$\% of edges with highest $U_{ij}$.
\item \textbf{Threshold:} remove edges with $U_{ij}>T$.
\end{itemize}
The result is $\bar{G}_i=(V,\bar{E}_i,D,Y)$ as in Eq.~\eqref{eq:cosine_prune}.
As with Method~II, each LLM view $i$ can induce a different $\bar{E}_i$; GNNs and
the DIVERGE ensemble proceed unchanged on the refined graphs.


% ---------------------------------------------------------------------------
\subsection{Overall Pipeline and Complexity}
\label{sec:pipeline-summary}
% ---------------------------------------------------------------------------

\textbf{Pipeline.} (1)~$M$ LoRA fine-tunes on native text $\rightarrow$
$\mathbf{Z}^{(i)}$; (2)~optional structure refinement
$\bar{G}_i$ via Method~II or III; (3)~$M$ GNNs on
$(\bar{G}_i,\mathbf{Z}^{(i)})$; (4)~adaptive ensemble over logits.

\textbf{Design choices.} Methods~II and III are modular and mutually exclusive per
experiment (or combinable in future work). Neither calls external LLM APIs for
structure editing. Method~II is simpler and directly encodes semantic homophily;
Method~III is model-driven and better suited when relational ambiguity dominates
\citep{yu2020grcn,zhao2021gaug}.

\textbf{Cost.} Over DIVERGE, Method~II adds $O(|E|)$ similarity checks per view;
Method~III adds $M_{\mathrm{cv}}$ edge-predictor training runs---typically small
relative to $M$ LoRA fine-tunes.

%=============================================================================
\section{Experiments}
\label{sec:experiments}
%=============================================================================

\subsection{Experimental Setup}
\textbf{Datasets.} We evaluate on four standard TAG benchmarks: Cora, CiteSeer, PubMed,
and WikiCS \citep{cora,giles1998citeseer,pubmed,mernyei2020wikics}. Each dataset
provides a citation or web-derived graph with text attributes and node-level class labels.

\textbf{Backbone.} We follow the DIVERGE pipeline: a pre-trained LLM is fine-tuned with
LoRA under $M$ random initializations (PiSSA, Orthogonal, EVA, Gaussian, LoFTQ),
producing $M$ semantic views. Each view feeds a GCN-style GNN trained on the (optionally
refined) graph. Three adaptive ensemble variants (Ensemble~1--3) combine per-view logits
with learnable weights.

\textbf{Baselines.} We compare against GCN, GraphSAGE, and GAT
\citep{Kipf2016GCN,hamilton2017sage,Velickovic2018GAT}; SenBERT and RoBERTa text
encoders; ENGINE; TAPE \citep{he2024harnessing}; GraphGPT \citep{tang2024graphgpt};
and LLaGA \citep{Chen2024LLaGA}.

\textbf{Structural enhancers.} Method~II applies cosine-similarity threshold pruning on
LLM embeddings (Sec.~\ref{sec:cosine-enhancer}). Method~III applies uncertainty-based
pruning using standard deviation or entropy from a cross-validated edge predictor
(Sec.~\ref{sec:uncertainty-enhancer}). Random removal of 30\% of edges serves as a
negative control.

\textbf{Metrics.} Node classification accuracy and Macro-F1 on held-out test nodes.


\subsection{Main Results}
Table~\ref{tab:diverge_with_cosine_enhancer_ensembles_accuracy} compares DIVERGE with
cosine-based semantic structure enhancement against strong baselines. Our method achieves
the best accuracy on all four datasets, reaching 94.10\% on Cora (vs.\ 90.05\% for TAPE)
and 97.13\% on PubMed. Table~\ref{tab:diverge_with_cosine_enhancer_ensembles_accuracy_macro_f1}
reports Macro-F1, where gains are similarly consistent.

\begin{table}[htbp]
\centering
\caption{Node classification accuracy (\%): DIVERGE + cosine enhancer vs.\ baselines.}
\label{tab:diverge_with_cosine_enhancer_ensembles_accuracy}
\small
\begin{tabularx}{\textwidth}{l|XXXX}
\toprule
\diagbox{\textbf{Dataset (\%)}}{\textbf{Method}} & Cora & CiteSeer & PubMed & WikiCS \\
\midrule
GCN & 87.41 & 75.74 & 89.01 & 83.67 \\
GraphSAGE & 87.44 & 74.96 & 90.47 & 84.86 \\
GAT & 86.68 & 73.73 & 88.25 & 83.94 \\
SenBERT-66M & 79.61 & 74.06 & 94.47 & 86.51 \\
RoBERTa-355M & 83.17 & 75.90 & 94.84 & 87.47 \\
ENGINE & 87.00 & 75.82 & 90.08 & 85.44 \\
TAPE & 90.05 & 76.45 & 93.00 & 87.11 \\
GraphGPT & 82.29 & 74.67 & 93.54 & 82.54 \\
LLaGA & 87.55 & 76.73 & 90.28 & 84.03 \\
\midrule
DIVERGE 1 + Enhancer & \textbf{93.91} & \textbf{81.03} & \textbf{97.03} & \textbf{89.32} \\
DIVERGE 2 + Enhancer & \textbf{94.10} & \textbf{81.82} & \textbf{97.11} & \textbf{89.19} \\
DIVERGE 3 + Enhancer & \textbf{94.10} & \textbf{81.50} & \textbf{97.13} & \textbf{89.28} \\
\bottomrule
\end{tabularx}
\end{table}

\begin{table}[htbp]
\centering
\caption{Macro-F1 (\%): DIVERGE + cosine enhancer vs.\ baselines.}
\label{tab:diverge_with_cosine_enhancer_ensembles_accuracy_macro_f1}
\small
\begin{tabularx}{\textwidth}{l|XXXX}
\toprule
\diagbox{\textbf{Dataset (\%)}}{\textbf{Method}} & Cora & CiteSeer & PubMed & WikiCS \\
\midrule
GCN & 86.54 & 71.52 & 88.54 & 82.11 \\
GraphSAGE & 86.37 & 71.87 & 90.16 & 82.78 \\
GAT & 85.64 & 69.27 & 87.70 & 82.14 \\
TAPE & 87.21 & 73.33 & 92.39 & 86.03 \\
LLaGA & 84.97 & 72.59 & 90.00 & 82.37 \\
\midrule
DIVERGE 1 + Enhancer & \textbf{92.65} & \textbf{77.92} & \textbf{96.60} & \textbf{88.43} \\
DIVERGE 2 + Enhancer & \textbf{92.99} & \textbf{78.56} & \textbf{96.74} & \textbf{88.27} \\
DIVERGE 3 + Enhancer & \textbf{92.96} & \textbf{78.37} & \textbf{96.77} & \textbf{88.09} \\
\bottomrule
\end{tabularx}
\end{table}

The results show that DIVERGE already achieves strong performance without refinement.
With cosine-based structural enhancement, further improvements are observed on most
datasets, indicating that graph structural enhancement acts as a complementary mechanism
for the DIVERGE framework.


\subsection{Cosine Similarity Analysis (Method II)}
To examine alignment between LLM representations and graph structure, we compute cosine
similarity between endpoint embeddings for same-label and different-label edges.
Table~\ref{tab:mean_std_cos_similarity} reports mean and standard deviation across LoRA
initialization schemes on Cora.

\begin{table}[htbp]
\centering
\caption{Mean and std.\ of cosine similarity for same-label vs.\ different-label edges (Cora).}
\label{tab:mean_std_cos_similarity}
\small
\begin{tabularx}{\textwidth}{l l *{5}{X}}
\toprule
Metric & Group & PiSSA & Orthogonal & EVA & Gaussian & LoFTQ \\
\midrule
Mean & same\_label & 0.818 & 0.903 & 0.810 & 0.895 & 0.799 \\
Mean & diff\_label & 0.729 & 0.851 & 0.724 & 0.842 & 0.723 \\
Std & same\_label & 0.099 & 0.059 & 0.097 & 0.063 & 0.092 \\
Std & diff\_label & 0.112 & 0.080 & 0.102 & 0.071 & 0.096 \\
\bottomrule
\end{tabularx}
\end{table}

\begin{figure}[htbp]
\centering
\maybeincludegraphics[width=0.95\textwidth]{Images/edge_similarity_boxplots_cora.png}
\caption{Cosine similarity distributions for same-label vs.\ different-label edges (Cora).}
\label{fig:cosine_box_plot}
\end{figure}

In all methods, the average cosine similarity for same-label edges is significantly
higher than for different-label edges, confirming that LLM representations align with
homophily. The higher standard deviation among different-label edges indicates structural
noise: some endpoints remain semantically close despite label mismatch. Figure~\ref{fig:cosine_box_plot}
shows clear distributional separation, supporting threshold-based pruning.

\subsection{Threshold Sensitivity}
\begin{figure}[htbp]
\centering
\maybeincludegraphics[width=0.75\textwidth]{Images/threshold_analysis_eva_cora.png}
\caption{Accuracy and edge-removal rate vs.\ cosine threshold (representative LoRA init on Cora).}
\label{fig:multi}
\end{figure}

Sensitivity analysis (Figure~\ref{fig:multi}) reveals a practical sweet spot: at low
thresholds almost no edges are removed; in a mid-range interval, modest pruning often
\emph{increases} accuracy; aggressive pruning eventually deletes useful bridges and
hurts performance.


\subsection{Comparison of Cosine and Uncertainty Enhancers}
Table~\ref{tab:accuracy_ensemble_after_enhancers} and
Table~\ref{tab:macro_f1_ensemble_after_enhancers} compare DIVERGE with and without
structural refinement across three ensemble variants.

\begin{table}[htbp]
\centering
\caption{Accuracy (\%): DIVERGE vs.\ cosine vs.\ uncertainty enhancers.}
\label{tab:accuracy_ensemble_after_enhancers}
\small
\begin{tabularx}{\textwidth}{l l *{3}{X}}
\toprule
Dataset & Approach & Ens.\ 1 & Ens.\ 2 & Ens.\ 3 \\
\midrule
\multirow{4}{*}{Cora}
 & DIVERGE & 93.17 & 93.36 & 93.17 \\
 & + Cosine & \textbf{93.91} & \textbf{94.10} & \textbf{94.10} \\
 & + Uncertainty (std) & 93.17 & 93.17 & 93.17 \\
 & + Uncertainty (entropy) & 92.99 & 92.80 & 92.62 \\
\midrule
\multirow{4}{*}{PubMed}
 & DIVERGE & 96.98 & 97.13 & 97.16 \\
 & + Cosine & 97.03 & 97.11 & 97.13 \\
 & + Uncertainty (std) & \textbf{97.26} & \textbf{97.16} & \textbf{97.16} \\
 & + Uncertainty (entropy) & 97.13 & 97.11 & 97.13 \\
\midrule
\multirow{4}{*}{CiteSeer}
 & DIVERGE & 81.66 & 81.97 & 82.13 \\
 & + Cosine & 81.03 & 81.82 & 81.50 \\
 & + Uncertainty (std) & 81.82 & 82.29 & 81.66 \\
 & + Uncertainty (entropy) & \textbf{83.07} & \textbf{82.76} & \textbf{82.92} \\
\midrule
\multirow{4}{*}{WikiCS}
 & DIVERGE & 88.34 & 88.38 & 88.42 \\
 & + Cosine & \textbf{89.32} & \textbf{89.19} & \textbf{89.28} \\
 & + Uncertainty (std) & 88.30 & 88.38 & 88.59 \\
 & + Uncertainty (entropy) & 88.64 & 88.51 & 88.64 \\
\bottomrule
\end{tabularx}
\end{table}

\begin{table}[htbp]
\centering
\caption{Macro-F1 (\%): DIVERGE vs.\ enhancers.}
\label{tab:macro_f1_ensemble_after_enhancers}
\small
\begin{tabularx}{\textwidth}{l l *{3}{X}}
\toprule
Dataset & Approach & Ens.\ 1 & Ens.\ 2 & Ens.\ 3 \\
\midrule
\multirow{4}{*}{Cora}
 & DIVERGE & 92.06 & 91.80 & 91.85 \\
 & + Cosine & \textbf{92.65} & \textbf{92.99} & \textbf{92.96} \\
 & + Uncertainty (std) & 91.96 & 91.96 & 91.96 \\
 & + Uncertainty (entropy) & 91.77 & 91.43 & 91.59 \\
\midrule
\multirow{4}{*}{CiteSeer}
 & DIVERGE & 79.25 & 78.83 & 79.04 \\
 & + Cosine & 77.92 & 78.56 & 78.37 \\
 & + Uncertainty (std) & 79.25 & 79.25 & 79.25 \\
 & + Uncertainty (entropy) & \textbf{80.72} & \textbf{80.72} & \textbf{80.72} \\
\bottomrule
\end{tabularx}
\end{table}

Cosine pruning peaks on Cora and WikiCS; uncertainty (std) helps on PubMed; entropy
helps on CiteSeer---confirming complementarity between Methods~II and~III.


\subsection{Random Edge Removal Baseline}
Table~\ref{tab:compare_structural_enhancer_random_model} compares targeted enhancers
against random 30\% edge deletion on Cora (representative LoRA inits). Random removal
causes large accuracy drops, whereas cosine and uncertainty refinement improve or
maintain performance.

\begin{table}[htbp]
\centering
\caption{Single-view accuracy (\%) on Cora: base vs.\ enhancers vs.\ random removal.}
\label{tab:compare_structural_enhancer_random_model}
\small
\begin{tabularx}{\textwidth}{l *{5}{X}}
\toprule
Approach & PiSSA & Orthogonal & EVA & Gaussian & LoFTQ \\
\midrule
Base & \textbf{90.77} & 89.48 & 92.07 & 91.14 & 90.96 \\
Cosine & \textbf{90.77} & 90.41 & 92.07 & 90.04 & \textbf{91.88} \\
Uncertainty (std) & 89.67 & 91.51 & \textbf{92.44} & 91.88 & 91.14 \\
Uncertainty (entropy) & 89.48 & \textbf{91.70} & \textbf{92.44} & \textbf{92.25} & 91.51 \\
Random 30\% & 88.56 & 89.59 & 90.88 & 90.14 & 89.41 \\
\bottomrule
\end{tabularx}
\end{table}


\subsection{Uncertainty Analysis (Method III)}
\begin{figure}[htbp]
\centering
\maybeincludegraphics[width=0.55\textwidth]{Images/uncertainty_boxplot_seaborn_cora.png}
\caption{Edge uncertainty distribution on Cora (entropy vs.\ std).}
\label{fig:edge_similarity_boxplots_cora}
\end{figure}

\begin{figure}[htbp]
\centering
\maybeincludegraphics[width=0.45\textwidth]{Images/std_uncertainty_vs_K_cora.png}
\hfill
\maybeincludegraphics[width=0.45\textwidth]{Images/entropy_uncertainty_vs_K_cora.png}
\caption{Effect of number of edge-prediction runs on average uncertainty (Cora).}
\label{fig:std_uncertainty_vs_K_cora}
\end{figure}

On Cora, entropy is more dispersed than standard deviation across edges. Increasing
the number of cross-validation runs stabilizes estimates with diminishing returns.
Std-based removal yields stable PubMed gains; entropy excels on noisier CiteSeer graphs.


%=============================================================================
\section{Discussion}
\label{sec:discussion}
%=============================================================================

In this section we interpret the experimental findings, relate them to prior work on
graph structure learning and LLM-based TAG methods, and discuss limitations and
future directions.


\subsection{Semantic Structure Enhancement as a Complement to DIVERGE}
\label{sec:disc-diverge-complement}

Our results support the central hypothesis of this work: \emph{rich LLM semantics and
clean graph topology are both necessary} for strong TAG node classification. DIVERGE
already achieves competitive or state-of-the-art performance by encoding native node
text through multiple LoRA-adapted LLMs and combining per-view GNNs with adaptive
ensembling rather than uniform logit averaging \citep{he2024harnessing}. Tables
\ref{tab:diverge_with_cosine_enhancer_ensembles_accuracy} and
\ref{tab:accuracy_ensemble_after_enhancers} show that further gains are available when
edges connecting semantically dissimilar nodes are removed before message passing.

This pattern aligns with graph structure learning (GSL) literature, which has repeatedly
shown that GNNs are sensitive to noisy or task-misaligned edges
\citep{jin2020prognn,chen2020idgl,dong2023gps}. What differs in our setting is the
\emph{source} of the feature signal used for denoising. Classical GSL methods such as
IDGL and ProGNN refine adjacency from a single shallow or jointly learned embedding
space \citep{chen2020idgl,jin2020prognn}. DIVERGE instead supplies multiple semantic
views from independently initialized LoRA fine-tunes; each view induces its own refined
graph $\bar{G}_i$, so structural enhancement amplifies representation diversity rather
than collapsing structure learning into one end-to-end optimization. The final ensemble
thus exploits three complementary axes: \textbf{semantic diversity} (multiple LLM
encodings), \textbf{structural diversity} (view-specific denoised topologies), and
\textbf{predictive diversity} (learned ensemble weights over GNN logits).

Compared with LLM \emph{enhancer} pipelines that rewrite node text, synthesize labels,
or call external models to modify topology
\citep{he2024harnessing,Sun2026TopologicalEnhancerTAG,guo2024graphedit,zhang2025leveraging},
our structural enhancer is deliberately lightweight: it requires no additional LLM API
queries at refinement time and preserves the native document content that DIVERGE was
designed to exploit. On Cora, CiteSeer, PubMed, and WikiCS, DIVERGE with cosine-based
structural refinement outperforms TAPE, LLaGA, and standard GNN baselines
(Tables~\ref{tab:diverge_with_cosine_enhancer_ensembles_accuracy}--\ref{tab:diverge_with_cosine_enhancer_ensembles_accuracy_macro_f1}),
suggesting that homophily-oriented edge pruning can recover part of the benefit otherwise
sought through generated explanations or rewritten inputs---while remaining compatible
with DIVERGE's no-augmentation philosophy.


\subsection{Why Cosine Similarity Identifies Noisy Edges}
\label{sec:disc-cosine}

The cosine-similarity enhancer (Method~II) rests on two empirical observations. First,
same-label edges exhibit systematically higher mean cosine similarity than
different-label edges in LLM embedding space (Table~\ref{tab:mean_std_cos_similarity},
Figure~\ref{fig:cosine_box_plot}), confirming that fine-tuned LLM representations align
with homophily to a useful degree. Second, different-label edges show greater variance,
including cases where endpoints remain semantically close despite label mismatch---a
signature of structural noise or graph-construction rules that ignore semantic class
boundaries.

Threshold-based pruning removes edges whose endpoints are far apart in representation
space before GNN propagation. Sensitivity analysis (Figure~\ref{fig:multi}) reveals a
practical sweet spot: at low thresholds almost no edges are removed and accuracy is
unchanged; in a mid-range interval, a modest fraction of edges is pruned and accuracy
often \emph{increases}, indicating that removed connections were harmful rather than
informative. Aggressive pruning eventually hurts performance by deleting useful
bridges---especially when heterophily or long-range dependencies matter
\citep{dong2023gps}. This trade-off parallels PTDNet and GPS, where feature-aware edge
removal helps only when sparsification is controlled \citep{luo2021ptdnet,dong2023gps}.

Random removal of 30\% of edges
(Table~\ref{tab:compare_structural_enhancer_random_model}) causes substantial accuracy
drops, whereas targeted cosine pruning improves or maintains performance. The contrast
confirms that gains are not an artifact of sparsifying the graph per se, but of removing
edges specifically inconsistent with LLM-estimated semantic proximity. On Cora, CiteSeer,
and WikiCS, cosine refinement yields the largest ensemble gains
(Table~\ref{tab:accuracy_ensemble_after_enhancers}), whereas on PubMed the improvement
is smaller---likely because citation links are already relatively homogeneous and the
baseline DIVERGE accuracy is already high ($\sim$97\%).


\subsection{Uncertainty-Based Refinement and Edge-Prediction GSL}
\label{sec:disc-uncertainty}

The uncertainty-based enhancer (Method~III) addresses a complementary failure mode.
Cosine similarity measures \emph{static semantic distance} in embedding space; it does
not explicitly ask whether an edge is supported by relational patterns learnable from
the rest of the graph. By training an edge predictor under cross-validation and flagging
edges with high standard deviation or entropy across folds, we target connections whose
existence is \emph{ambiguous}---a criterion closer to edge-prediction GSL such as GRCN
and GAug \citep{yu2020grcn,zhao2021gaug}.

On Cora (Figure~\ref{fig:edge_similarity_boxplots_cora}), entropy is more sensitive and
dispersed than standard deviation. Accordingly, std-based removal tends to yield stable
improvements on PubMed (Table~\ref{tab:accuracy_ensemble_after_enhancers}), while
entropy-based removal can excel on CiteSeer. Increasing the number of edge-prediction
runs stabilizes uncertainty estimates (Figure~\ref{fig:std_uncertainty_vs_K_cora}) but
with diminishing returns beyond a modest fold count.

Neither uncertainty measure dominates uniformly. Cosine pruning often achieves larger
peak accuracy (e.g., Cora 94.10\% vs.\ 93.36\% baseline), while uncertainty pruning
can improve Macro-F1 on imbalanced or noisier graphs (e.g., CiteSeer with entropy;
PubMed with std in Table~\ref{tab:macro_f1_ensemble_after_enhancers}). This
complementarity suggests that semantic distance and relational ambiguity capture
partially overlapping but non-identical sets of noisy edges.


\subsection{Choosing Between Enhancers in Practice}
\label{sec:disc-when-to-use}

Based on our experiments, we offer the following practical guidance:

\begin{itemize}
\item \textbf{Cosine similarity (Method~II)} is preferred when LLM embeddings show
  clear separation between same-label and different-label edges
  (Table~\ref{tab:mean_std_cos_similarity}), when homophily is a reasonable approximation
  of the downstream task, and when interpretability matters. It delivers the strongest
  gains on Cora and WikiCS in our evaluation.

\item \textbf{Uncertainty with standard deviation (Method~III, std)} is preferred when
  the graph is large, already fairly homophilic, and incremental but stable improvement
  is desired (PubMed).

\item \textbf{Uncertainty with entropy (Method~III, entropy)} is worth considering on
  graphs where label noise or heterophily produces many ambiguous links (CiteSeer).

\item \textbf{DIVERGE multi-view vs.\ single-view ensemble:} Single-representation
  ensemble on multiple refined graphs approaches full DIVERGE performance on most
  datasets, offering a cost-effective deployment option when training multiple LoRA
  adapters is prohibitive.
\end{itemize}

A natural extension is to combine criteria, e.g., removing edges that simultaneously
fall below a cosine threshold \emph{and} rank among the highest-uncertainty links.


\subsection{Relation to Graph Structure Learning and LLM Enhancers}
\label{sec:disc-prior-work}

Our semantic structure enhancer occupies a middle ground between GNN-only GSL and
LLM-based TAG enhancers (Section~\ref{sec:related-works}). Unlike ProGNN, IDGL, and
PTDNet, we do not jointly learn a single adjacency matrix inside one training loop;
refinement is a modular post-processing step on each LLM view. Unlike RoSE, GraphEdit, and
LLM4GraphTopology \citep{Sun2026TopologicalEnhancerTAG,guo2024graphedit}, we do not
invoke LLMs to rewrite edges or decompose relation types; structure is inferred from
embeddings and GNN edge predictions already computed for classification.

This design trades the expressive power of generative LLM structure editing for
\emph{efficiency and reproducibility}: no external API calls, no hallucinated edge
labels, and transparent criteria (cosine or uncertainty) for each removed link. The cost
is that we cannot add missing edges or discover latent relation types that never appear
in the observed graph---capabilities that LLM enhancers explicitly target
\citep{Sun2026TopologicalEnhancerTAG,Yu2025NodeGenerationTAG}. For many citation and
web-derived TAG benchmarks where the primary issue is \emph{spurious} rather than
\emph{missing} edges, pruning-oriented enhancement appears sufficient.


\subsection{Limitations}
\label{sec:disc-limitations}

Several limitations should be acknowledged.

\paragraph{Homophily assumption.}
Both enhancers implicitly assume that removing low-similarity or high-uncertainty edges
helps classification. On strongly heterophilic graphs, beneficial edges may connect
dissimilar nodes; methods such as ES-GNN explicitly split rather than delete such
links \citep{li2024esgnn}. Our evaluation focuses on citation-style benchmarks where
homophily-oriented refinement is plausible but not guaranteed.

\paragraph{Hyperparameter sensitivity.}
Cosine threshold $T$, uncertainty threshold, and top-$k$ removal rate must be tuned
per dataset and often per LLM view (Figure~\ref{fig:multi}). Unlike end-to-end GSL,
refinement is not automatically aligned with the downstream loss.

\paragraph{Computational overhead.}
Method~II adds pairwise cosine computation over edges for each of $M$ LLM views;
Method~III adds cross-validated edge-predictor training. Both costs are typically
smaller than multiple LoRA fine-tunes but non-negligible on million-edge graphs.

\paragraph{Scope of evaluation.}
We study node classification on four TAG benchmarks. Extension to link prediction,
graph-level tasks, dynamic graphs, or graphs with rich edge text remains open. We also
do not compare against every recent LLM enhancer (e.g., UltraTAG, SKETCH, TAGA)
\citep{yu2025ultratag,zhou-etal-2025sketch,zhang2025taga} under identical tuning budgets.

\paragraph{CiteSeer Macro-F1.}
Cosine refinement improves CiteSeer accuracy but can reduce Macro-F1 relative to
baseline DIVERGE (Table~\ref{tab:macro_f1_ensemble_after_enhancers}), suggesting that
aggressive pruning may hurt minority classes.


\subsection{Future Work}
\label{sec:disc-future}

Promising directions include:
(i)~ \textbf{adaptive, per-class or per-view thresholds} learned from validation data;
(ii)~ \textbf{hybrid enhancers} intersecting cosine and uncertainty scores, or combining
pruning with edge addition as in GAug \citep{zhao2021gaug};
(iii)~ \textbf{integration with LLM enhancers} in low-data regimes
\citep{zhang2025leveraging,hu2025gaga};
(iv)~ \textbf{theoretical analysis} connecting removed-edge rate to homophily ratio,
building on ESNR and minimal-sufficient-structure frameworks
\citep{dong2023gps,fatemi2022cogsl};
(v)~ \textbf{large-scale and heterophilic benchmarks} (ogbn-products, Roman-Empire);
(vi)~ \textbf{single-pass deployment} combining the best refined graphs from one LLM
view with multi-graph ensemble learning.


%=============================================================================
\section{Conclusion}
\label{sec:conclusion}
%=============================================================================

We extended the DIVERGE framework for text-attributed graph node classification with a
\emph{semantic structure enhancer} that removes noisy edges before GNN message passing,
without rewriting native node text or calling external LLMs for topology editing. Two
mechanisms---cosine-similarity thresholding on multi-view LLM embeddings and
uncertainty-based pruning via cross-validated edge prediction---increase effective
homophily and improve ensemble accuracy over DIVERGE alone on Cora, CiteSeer, PubMed,
and WikiCS. Targeted removal consistently outperforms random edge deletion, confirming
that structure quality mediates the benefit of rich semantic representations. Cosine and
uncertainty criteria are complementary: the former excels when embedding-space homophily
is strong; the latter offers stable gains when relational ambiguity dominates. By
combining native-text LLM diversity with GSL-inspired denoising, our approach offers an
efficient alternative to full LLM enhancer pipelines while pushing TAG performance
beyond strong baselines including TAPE. Future work will pursue adaptive hybrid
enhancers, broader benchmarks, and tighter integration with cost-aware LLM--GNN systems.


\bibliographystyle{plainnat}
\bibliography{references,sample_extra}

\end{document}
